\section{Second Fundamental Theorem of Calculus (FTC2)} 

\subsection{FTC2}
FTC2:

Given a continuous function $f(x)$ IF 
   \[
	G(x) = \int_{a}^{x} f(t) \; dt 
   \]
   Where $t$ is dummy variable and varies between $a$ and $x$ 

   THEN,

      \[
   	\dot{G}(x) = f(x)
      \]
      Geometrically , 
\begin{figure}[H]
	\centering
	\def\svgwidth{0.7\columnwidth}
	\import{./images/}{FTC2.pdf_tex}
\end{figure}
$G(x)$ varies as $x$ varies , $G(x)$ is a function itself and gives 
area between $y =  f(t)$ and $t$ axis and between lines $t=a$ and $t=x$

In terms of differential equation , 
FTC2 says that $G(x)$ is a solution to 
   \begin{align*}
	\dot{y} = f \; \text{differential equation} \\
	y(a) = 0 \; \text{initial equation} 
   \end{align*}
   $y(a)= 0 $ follows from 
      \[
   	G(a) = \int_{a}^{a} f(t) \; dt = 0
      \]
The integrand $f(x)$ can be any continous function and Not  the just the 
ones whose antiderivative we know how to find. These integral formula are still usefull ,
There are numerical methods that allows us to compute them

\subsection{FTC2 and Chain Rule}
for any continous function $f$ and differential function $u(x)$ 
\begin{align*}
    &\frac{d}{dx} \int_{a}^{u(x)} f(t) \; dt = f(u(x)) \times \dot{u}(x) \\
    =& \frac{d}{d(u(x))} \int_{a}^{u(x)} f(t) \; dt \times \frac{d(u(x))}{d(x)}\\
    =& f(u(x)) \times \dot{u}(x)
\end{align*}

\subsection{The two fundamental Theorem of Calculus}
\begin{tabular}{l|l}
	\toprule
	FTC 1 & FTC2 \\
	\midrule
	Given a differential function F with continuous derivative $\dot{F}$ &               Given a continuous function $f$ \\
	$\int_{z}^{x} \dot{F} \; dt = F(x) - F(x)$                           & $\frac{d}{dx} \int_{z}^{x} f(t) \; dt = f(x)$\\
	\bottomrule
\end{tabular}

The two fundamental theorem together say that differentiation and integration
(from $a$ to $x$)  are inverse operation of one another up 
to a constant. 

\subsection{Constructing Function using Integral}
Recall that, FTC2 gives the integral formula 
   \[
	G(x) = \int_{a}^{x} f(t) \; dt
   \]
   as solution to following IVP
      \[
   	\dot{y} = f \; y(a) = 0
      \]
Depending on what differential equation is $G(x)$ can be a function we 
already know or the one that cannot be expressed (without and integral) 
in terms of any function we already know . 

Ex,The logarithmic function is defined using integral 
\begin{enumerate}

    \item 
   \[
	L(x) = \int_{1}^{x} \frac{dt}{t} = ln(x)
   \]
   \item  The integral of $e^{-t^{2}}$ is called the bell curve . it is a function 
	defined by integral and cannot be expressed in  terms of function we 
	already know. 
	   \[
		F(x) = \int_{0}^{x} e^{-t^{2}} 
	   \]
\end{enumerate}
\paragraph{Definition of Logarithm and exponential  function}
exponential function $E(y)$ is solution to following Initial Value problem .
   \[
	\dot{E}(y) = E(y) \; E(0) = 1
   \]
   We then let $e = E(1)$ and we denote $E(y)$ by $e^{y}$

   Logarithmic function : Inverse of exponential function $E^{-1}(x)$, denoted by $ln(x)$
   i.e 
      \[
   	E^{-1}(x) = ln(x)
      \]
\subsection{The integration of bell curve }

bell curve is the graph of $e^{-t^{2}}$ 
   \[
	F(x)  = \int_{0}^{x} =  e^{-t^{2}} \; dt
   \]
   Geometrically $F(x)$ denote the following shaded Region. 
\begin{figure}[H]
	\centering
	\def\svgwidth{0.5\columnwidth}
	\import{./images/}{bellcurve.pdf_tex}
\end{figure}

using First and Second derivative of $F$ we sketch the graph of $F$ as 
shown below 

\begin{figure}[H]
	\centering
	\def\svgwidth{0.5\columnwidth}
	\import{./images/}{area_bell_curve.pdf_tex}
\end{figure}
Two properties of $F$
\begin{enumerate}

    \item $F$ is odd $F(-x) = - F(x) $
	\item $\lim_{x \to \infty} F(x) = \frac{\sqrt{\pi}}{2}$ and $\lim_{x \to \infty} F(x) = - \frac{\sqrt{\pi}}{2}$

\end{enumerate}
Geometrically, $\lim_{x \to  \infty}  F(x)$  is the area of shaded infinite sketch. 
\begin{figure}[H]
	\centering
	\def\svgwidth{0.5\columnwidth}
	\import{./images/}{bellcurve1.pdf_tex}
\end{figure}

\subsection{ The error Function}
The error function denoted by $erf(x)$ is renormalized version of $F(x) = \int_{0}^{x} e^{-t^{2}} \; dt$ 
that has horizental asymptotes at $y \pm 1$ 
   \[
	erf(x) = \frac{2}{\sqrt{pi}} F(x) = \frac{2}{\sqrt{pi}} \int_{0}^{x} e^{-t^{2}} \; dt
   \]
   \subsubsection{More new function} 
   TODO	 from Bookmarked

